\chapter{Solid State Theory: Midterm, 5 December 2014}
\label{q:theory-2014-12-midterm}
\begin{sourceNote}
Both printed pages survive. The printed question weights sum to 110\%; they are retained without renormalization. Handwritten answers and annotations are not part of the transcribed questions.
\end{sourceNote}

\section*{1. (40\%) Explain the following terms.}
\begin{enumerate}[label=(\alph*)]
\item Bravais lattices;
\item composite crystals;
\item Brillouin zone;
\item reciprocal lattices;
\item Bloch theorem;
\item exchange energy of a many-electron system;
\item fully spin-polarized ground state;
\item Wigner crystallization.
\end{enumerate}

\section*{2. (20\%) Lattice planes and periodicity}
Consider a Bravais lattice spanned by the primitive vectors
\[
\vct{t}_1=\frac a2(0,1,1),\qquad
\vct{t}_2=\frac a2(1,0,1),\qquad
\vct{t}_3=\frac a2(1,1,0).
\]
$\vct{g}_1$, $\vct{g}_2$, and $\vct{g}_3$ are the primitive vectors of the reciprocal lattice, defined by the relations $\vct{t}_i\cdot\vct{g}_j=2\pi\delta_{ij}$.
\begin{enumerate}[label=(\alph*)]
\item (5\%) Consider a family of direct lattice planes that is normal to $3\vct{g}_1+6\vct{g}_2-6\vct{g}_3$. Find the distance between two successive lattice planes.
\item (5\%) The lattice points on each lattice plane in (a) form a two-dimensional lattice. Find the density of lattice points (i.e., the number of lattice points per unit area) on the two-dimensional lattice plane.
\item (5\%) Consider a family of direct lattice planes that is normal to $\frac12\vct{t}_1+\vct{t}_2-\vct{t}_3$. Find the distance between two successive lattice planes.
\item (5\%) Does the plane wave $\exp(\ii\frac12\vct{g}_1\cdot\vct{r})$ have the lattice periodicity? Explain for your answer.
\end{enumerate}

\section*{3. (20\%) Band critical points}
Consider a two-dimensional crystal that is underlain by a two-dimensional hexagonal lattice spanned by the primitive vectors
\[
\vct{t}_1=a\left(\frac12,\frac{\sqrt3}{2}\right),\qquad
\vct{t}_2=a\left(-\frac12,\frac{\sqrt3}{2}\right).
\]
The crystal has two bands with the dispersion relations
\[
E_{\vct{k}}=E_p\pm\gamma\abs{F_{\vct{k}}},\qquad \gamma>0,
\]
where
\[
\abs{F_{\vct{k}}}^{2}
=1+4\cos\frac{k_xa}{2}\cos\frac{\sqrt3 k_ya}{2}
+4\cos^2\frac{k_xa}{2}.
\]
\begin{enumerate}[label=(\alph*)]
\item (10\%) Identify if or not symmetry points $\Gamma$ ($\vct{k}=0$) and $P$ [$\vct{k}=(2\pi/a)(2/3,0)$] of the first Brillouin zone are critical points. What type of critical point are they if they are critical points?
\item (10\%) Express the density of states contributed from the states in a small region containing a critical point you choose as you like from those in (a).
\end{enumerate}

\section*{4. (5\%) Many-electron Hamiltonian}
Express the Hamiltonian for a system of $N$ interacting electrons in the presence of nuclei fixed in a configuration.

\section*{5. (5\%) Hartree trial state}
Describe the trial wavefunction for the ground state of a many-electron system in the Hartree theory.

\section*{6. (10\%) Hartree--Fock trial state}
Describe the trial wavefunction for the ground state of a many-electron system in the Hartree--Fock theory. How to obtain the approximate ground-state wavefunction?

\section*{7. (10\%) Monte Carlo trial state}
Describe the trial wavefunction in the Monte Carlo Method of Ceperley and Alder for the ground state of a many-electron system in a uniform neutralizing background of positive charges.
